\chapter{The Moduli with Primitive Roots}

We have established the existence of primitive roots modulo any prime. In this section, we will give a complete characterization of those moduli for which a primitive root exists.

\section{Primitive Roots Modulo Odd Prime Powers}
Our first target is to show that any odd prime power $p^k$ has a primitive root. Our approach to this problem is of an arithmetic nature, based on an analysis of the $p$-valuations of differences of powers.

Recall that for any integer $a$ not divisible by $p$, the integer $a^{p-1}-1$ is divisible by $p$ by \zcref{theorem:fermats-little-theorem}. Therefore, its $p$-adic valuation $\alpha:=\nu_p(a^{p-1}-1)$ is at least $1$. In the following lemma we will establish that the order of $a$ modulo $p^k$ can be determined in terms of $\alpha$.

\begin{lemma}\label{lemma:order-modulo-prime-power-primitive-root}
  Consider an odd prime $p$ and a primitive root $g$ modulo $p$. Furthermore, let
  \begin{equation*}
    m:=\nu_p(g^{p-1}-1).
  \end{equation*}
  Then the order of $g$ modulo $p^k$ is $p^{\max(0,k-m)}(p-1)$. 
\end{lemma}

\begin{proof}
  First, we observe that $\ord_p(g)\mid \ord_{p^k}(g)$, so the order of $g$ modulo $p^k$ must be a multiple of $p-1$. Write $g^{p-1}=kp^{m}+1$, so that $p\nmid k$. By \zcrefnamed{theorem:lte} it follows that
  \begin{equation*}
    \nu_p(g^{p^t(p-1)}-1)=m+t.
  \end{equation*}
  Consequently, for any integer $s$ not divisible by $p$ we have
  \begin{equation*}
    g^{p^ts(p-1)}=(kp^{m+t}+1)^s\equiv 1+skp^{m+t}\pmod{p^{m+t+1}}.
  \end{equation*}
  This shows that $\nu_p(g^{p^ts(p-1)}-1)=m+t$ for such $s$. The first exponent $r$ such that $g^r\equiv 1\pmod{p^k}$ is therefore $r=p^{\max(0,k-m)}(p-1)$.
\end{proof}

\begin{theorem}\label{theorem:primitive-roots-odd-prime-powers}
  Any prime power $p^k$ with $k\geq 1$ admits a primitive root.
\end{theorem}

\begin{proof}
  By \zcref{lemma:order-modulo-prime-power-primitive-root} it suffices to construct a primitive root $h$ modulo $p$ such that $\nu_p(h^{p-1}-1)=1$.

  Consider a primitive root $g$ modulo $p$. Write $g^{p-1}=kp+1$, and define $h_t:=g+tp$. Then we have
  \begin{equation*}
    h_t^{p-1}=(g+tp)^{p-1}\equiv g^{p-1}+(p-1)g^{p-2}tp\equiv 1+p(k-g^{p-2}t) \pmod{p^2}.
  \end{equation*}
  Thus we see that $h_t^{p-1}\equiv 1\pmod{p^2}$ if and only if
  \begin{equation*}
    k \equiv g^{p-2}t \pmod{p}.
  \end{equation*}
  This linear congruence has the unique solution $t\equiv kg\pmod{p}$. Thus we see that for any $t\nequiv kg\pmod{p}$, we have $h_t^{p-1}\nequiv 1\pmod{p^2}$ and hence
  \begin{equation*}
    \nu_p(h_t^{p-1}-1)=1.\qedhere
  \end{equation*}
\end{proof}

\section{The Complete Characterization of Moduli with Primitive Roots}

\begin{theorem}\label{theorem:moduli-with-primitive-roots}
  Consider a positive integer $n$. The following are equivalent:
  \begin{enumerate}
  \item \label{itemi:classification-primitive-roots}The number $n$ admits a primitive root.
  \item \label{itemii:classification-primitive-roots}Either $n=2$, or there is a unique element of order $2$ modulo $n$.
  \item \label{itemiii:classification-primitive-roots}We have
  \begin{equation*}
    n=2,~4,~p^k,~\text{or }~2p^k,
  \end{equation*}
  where $p$ is an odd prime.
  \end{enumerate}
\end{theorem}

\begin{proof}
  We saw in \zcref{theorem:primitive-roots} that if $n$ admits a primitive root, then the congruence
  \begin{equation*}
    x^d\equiv 1\pmod{n}
  \end{equation*}
  had at most $d$ solutions, for any $d\geq 0$. In particular, there are at most two solutions to the congruence $x^2\equiv 1\pmod{n}$. One of them is the number $1$, which is of order $1$. The other is $-1$, which is of order $2$. Since any element of order $2$ is a solution to the congruence $x^2\equiv 1\pmod{n}$, we see that \ref{itemi:classification-primitive-roots} implies \ref{itemii:classification-primitive-roots}.

  Next, assume that there is a unique element of order $2$ modulo $n$. The case $n=4$ is swiftly dealt with, since $-1$ is the unique integer of order $2$ modulo $4$. If $n=2^k$ for some $k\geq 3$, then there are three integers of order $2$ modulo $2^k$:
  \begin{equation*}
    x=- 1\qquad\text{and}\qquad x=2^{k-1}\pm 1.
  \end{equation*}
  Thus we see that if $n=2^k$ for some $k\geq 3$, then $n$ does not satisfy the hypothesis in \ref{itemii:classification-primitive-roots}.

  Now suppose that $n=km$, where $\gcd(k,m)=1$ and $m$ is odd and strictly greater than $1$. By \zcrefnamed{theorem:chinese-remainder-theorem}, it follows that any pair of solutions
  \begin{align*}
    x^2 & \equiv 1\pmod{k}, \\
    y^2 & \equiv 1\pmod{m},
  \end{align*}
  determines a unique solution to the congruence $z^2\equiv 1\pmod{km}$. It follows that one of them can have only one nontrivial solution. Since we assumed that $m>1$ is odd, and thus has two distinct solutions, it follows that there is only one solution to the congruence $x^2\equiv 1\pmod{k}$. Observe that there is only one such solution if and only if $k=1$ or $k=2$. Furthermore, the previous observation implies that $m$ cannot have distinct prime factors, so we conclude that $m$ must be of the form $m=p^k$ for some odd prime $p$. Thus we have proven that \ref{itemii:classification-primitive-roots} implies \ref{itemiii:classification-primitive-roots}.

  Finally, we have to show that any number of the form $n=2,4,p^k$, or $2p^k$, where $p$ is an odd prime, admits a primitive root. The number $1$ is a primitive root modulo $2$, and the number $-1$ is a primitive root modulo $4$. The fact that any odd prime power $p^k$ admits primitive roots was established in \zcref{theorem:primitive-roots-odd-prime-powers}. To finish the proof that \ref{itemiii:classification-primitive-roots} implies \ref{itemi:classification-primitive-roots}, it therefore remains to show that each number of the form $n=2p^k$, where $p$ is an odd prime, has a primitive root. To see this, observe that the integers relatively prime to $2p^k$ are exactly the odd integers $1\leq x<2p^k$ relatively prime to $p^k$. Among those odd integers is a primitive root $g$ modulo $p^k$. Its order modulo $2p^k$ is
  \begin{equation*}
    \ord_{2p^k}(g)=\ord_{p^k}(g)=p^{k-1}(p-1)=\phi(2p^k),
  \end{equation*}
  so we see that $g$ is also a primitive root of $2p^k$.
\end{proof}

\section{Carmichael's Function}

%%% Local Variables:
%%% mode: latex
%%% TeX-master: "number-theory"
%%% TeX-engine: luatex
%%% End:
